\section{Einleitung / Problemstellung}

\begin{frame}{Auffrischer Clustering}
    \begin{minipage}[t]{0.6\textwidth}
        \begin{itemize}
            \item zur Analyse von Daten
            \begin{itemize}
                \item unsupervised
                \item Nähe von Datenpunkten finden
            \end{itemize}
            \item Zielfunktion hier k-center
            \begin{itemize}
                \item maximaler Radius minimieren
                \item Zentren aus der Punktemenge
            \end{itemize}
            \item optimales k-center $\in NP$ \footfullcite{Clustering_Skript}
            \item Algorithmus von Gonzalez als bekannter Vertreter
        \end{itemize}
    \end{minipage}
    \hfill
    \begin{minipage}[t]{0.35\textwidth}
        \vspace{-\baselineskip} % Move image to the top
        % Beispiel für ein einfaches Clustering
        \includegraphics[width=\textwidth]{example-image-a}
    \end{minipage}

\end{frame}











\begin{frame}{Gonzalez}
    \begin{minipage}[t]{0.6\textwidth}
        \begin{itemize}
            \item Appoximationsgarantie von 2
            \item colorblind
        \end{itemize}
        \begin{algorithm}[H]
            \scriptsize
            \begin{algorithmic}[1]
                % Anfangsbedingungen
                \State $Z := $ leere Zentrenmenge
                \State Wähle erstes Zentrum $z_0$ zufällig

                % Daten für erstes Zentrum aktuallisieren
                \ForAll{$x \in P$}
                    \State $d_{min} = dist(z_0, x)$
                \EndFor

                % Für Jedes K durchgen
                \For{$i \in 0,..,k-1$}
                    % bestes Zentum suchen
                    \State $z_i = P(argmax(d_{min}))$

                    % Alle Punkte aktuallisieren
                    \For{$x \in P \setminus Z$}
                        % Dieser Punkt muss aktuallisiert werden
                        \State Prüfe ob Punkte näher an neuem Zentrum
                    \EndFor
                \EndFor
                \State $r_{max} = $ größter Abstand zwischen Punkt und seinem Zentrum
                \State return $(Z, P, r_{max})$
            \end{algorithmic}
        \end{algorithm}
    \end{minipage}
    \hfill
    \begin{minipage}[t]{0.35\textwidth}
        \vspace{-\baselineskip} % Move image to the top
        % Beispiel für Gonzalez ausführung
        \includegraphics[width=\textwidth]{example-image-a}
    \end{minipage}
\end{frame}












\begin{frame}{Auffrischer fair Clustering}
    \begin{minipage}[t]{0.6\textwidth}
        \begin{itemize}
            \item faires Attribut
            \item Zusatzbedingung: gleichmäßig aufteilen auf Cluster
            \item optimales fair k-center ebenfalls $\in NP$ \footfullcite{Clustering_Skript}
        \end{itemize}
        % Unfaires Clustering mit Farben
        \includegraphics[width=\textwidth /2]{example-image-b}
    \end{minipage}
    \hfill
    \begin{minipage}[t]{0.35\textwidth}
        \vspace{-\baselineskip} % Move image to the top
        % Faires Clustering mit Farben
        \includegraphics[width=\textwidth]{example-image-a}
    \end{minipage}

\end{frame}


